\scanpage{93}
% archive foliation: 94
\figure{\begin{tikzpicture}\node[draw,ellipse]{$B$};\end{tikzpicture}}
Cas des groupes (commutatifs pour fixer les idées) \uline{plats}.

\margin{introduit par \unsure{Mazur-Robert}}
On pose
\diagram{\begin{array}{lll}
\ell_\bullet^{G/S}=\mathbf Le^*(L_\bullet^{G/S})&&\text{complexe de co-Lie}\\
\overset{\vee}{\ell}_\bullet^{G/S}=R\underline{\operatorname{Hom}}(\ell_\bullet^{G/S},\underline O_S)&&\rule{3cm}{0.3pt}\ \text{Lie}
\end{array}}
\margin{\uline{variances}}

Si $G$ est \uline{loc. de prés. finie} sur $S$, il est d'int. complèt rel. et on a
\diagram{\left\{\begin{array}{lll}
\ell_\bullet^{G/S}&\text{parfait d'ampl. parfaite}&\subset[-1,0]\\
\overset{\vee}{\ell}_\bullet^{G/S}&\rule{5cm}{0.3pt}&\subset[0,1]
\end{array}\right.}

Notations
\deleted{$\unsure{\underline\omega_G,\underline n_G}$}
\diagram{\left\{\begin{array}{l}
\underline\omega_G=\underline H_0(\ell_\bullet^{G/S})\\
\underline n_G=\underline H_1(\ell_\bullet^{G/S})\\
\underline t_G=H^0(\overset{\vee}{\ell}_\bullet^{G/S})\\
\nu_G=H^1(\overset{\vee}{\ell}_\bullet^{G/S})
\end{array}\right.\qquad t_G=\overset{\vee}{\omega}_G,\quad\underline n_G=\overset{\vee}{\nu}_G}
\[
\operatorname{rg}\ell_\bullet^G=\dim G\qquad(=\operatorname{rg}\underline\omega_G-\operatorname{rg}n_G\ \text{quand }\underline\omega_G,\underline n_G\text{ sont loc. libres}).
\]

Suite exacte de transitivité: Une suite exacte
\[
0\longrightarrow G'\longrightarrow G\longrightarrow G''\longrightarrow0
\]
de schémas en groupes plats implique une suite exacte
\diagram{\begin{tikzcd}&\ell_\bullet^{G'}\arrow[dl,"1"',very thick]&\\
\ell_\bullet^{G''}\arrow[rr]&&\ell_\bullet^G\arrow[ul]
\end{tikzcd}\quad\unsure{0}}
,
d'où, si $G'$ loc. de prés. finie, suite exacte courte
\[
0\longrightarrow n_{G'}-n_G\longrightarrow n_G\longrightarrow\omega_{G''}-\omega_G-\omega_{G'}-0.
\]

Formule \uline{remarquable}
\[
\boxed{R\underline{\operatorname{Hom}}(\ell_\bullet^G,J)\simeq\tau_{\leq1}R\operatorname{Hom}(G^{\unsure{*}},J)}
\]
\margin{\uuline{Faisant} $J=\underline O_S$, $G\mapsto G^*$, \uline{\unsure{à vérifier !}}
\diagram{\left\{\begin{array}{l}
\underline{\operatorname{Hom}}(G^*,\underline O_S)\simeq\underline t_{G^*}\\
\underline{\operatorname{Ext}}^1(G,\underline O_S)\simeq\underline\nu_G
\end{array}\right.}}
